\documentclass[12pt]{article}
\usepackage[UTF8]{inputenc}
\usepackage{thaisupp}
\usepackage{enumitem}
\begin{document}
\title{ไฟฟ้าสถิต}
\maketitle
\section*{In-class Questions}
\begin{enumerate}[label=\arabic*.]
\item ข้อใดต่อไปนี้ไม่ใช่หน่วยของประจุไฟฟ้า? \hfill \textbf{\small (Main)}
\begin{enumerate}[label=)]
\item ก) คูลอมบ์ (C)
\item ข) ไมโครคูลอมบ์ (μC)
\item ค) จูล (J)
\item ง) นาโนคูลอมบ์ (nC)
\end{enumerate}
\item ประจุไฟฟ้าบวกและประจุไฟฟ้าลบมีปฏิสัมพันธ์กันอย่างไร? \hfill \textbf{\small (Main)}
\begin{enumerate}[label=)]
\item ก) ขับกัน (ผลักกัน)
\item ข) ดึงกัน
\item ค) ไม่มีปฏิสัมพันธ์กัน
\item ง) ขึ้นอยู่กับระยะทาง
\end{enumerate}
\item กฎคูลอมบ์ระบุว่า แรงระหว่างประจุสองประจุแปรผันตรงกับข้อใด? \hfill \textbf{\small (Main)}
\begin{enumerate}[label=)]
\item ก) ผลคูณของขนาดประจุทั้งสอง และแปรผกผันกับระยะทางกำลังสอง
\item ข) ผลคูณของขนาดประจุทั้งสอง และแปรผกผันกับระยะทาง
\item ค) ผลรวมของขนาดประจุทั้งสอง และแปรผันตรงกับระยะทางกำลังสอง
\item ง) ผลคูณของขนาดประจุทั้งสอง และแปรผันตรงกับระยะทางกำลังสอง
\end{enumerate}
\item ประจุสองประจุ q₁ = +2 μC และ q₂ = -3 μC วางห่างกัน 3 cm ในสุญญากาศ ค่าคงตัว k = 9×10⁹ N·m²/C² จงหาขนาดของแรงที่กระทำต่อประจุแต่ละตัว \hfill \textbf{\small (Main)}
\begin{enumerate}[label=)]
\item ก) 3 N
\item ข) 6 N
\item ค) 9 N
\item ง) 60 N
\end{enumerate}
\item สนามไฟฟ้า ณ จุดใดๆ หมายถึงข้อใด? \hfill \textbf{\small (Main)}
\begin{enumerate}[label=)]
\item ก) แรงที่กระทำต่อประจุบวก 1 คูลอมบ์ ที่จุดนั้น
\item ข) พลังงานที่จุดนั้นต่อหน่วยประจุ
\item ค) ศักย์ไฟฟ้าที่จุดนั้น
\item ง) กระแสไฟฟ้าที่จุดนั้น
\end{enumerate}
\item จุดประจุ Q = +5 μC สร้างสนามไฟฟ้ารอบตัว ณ ตำแหน่งที่ห่างจาก Q เป็นระยะ 2 m สนามไฟฟ้ามีขนาดเท่าใด? (k = 9×10⁹ N·m²/C²) \hfill \textbf{\small (Main)}
\begin{enumerate}[label=)]
\item ก) 1.125×10⁴ N/C
\item ข) 2.25×10⁴ N/C
\item ค) 1.125×10³ N/C
\item ง) 2.25×10³ N/C
\end{enumerate}
\item ประจุบวกเคลื่อนที่ในสนามไฟฟ้าสม่ำเสมอจากจุด A ไปยังจุด B ที่อยู่ห่างกัน 0.5 m โดยทิศของสนามไฟฟ้าชี้จาก A ไป B ความต่างศักย์ระหว่าง A กับ B เป็น 100 V ขนาดของสนามไฟฟ้ามีค่าเท่าใด? \hfill \textbf{\small (Main)}
\begin{enumerate}[label=)]
\item ก) 50 V/m
\item ข) 100 V/m
\item ค) 150 V/m
\item ง) 200 V/m
\end{enumerate}
\item ศักย์ไฟฟ้า (Electric Potential) คือข้อใด? \hfill \textbf{\small (Main)}
\begin{enumerate}[label=)]
\item ก) งานที่ใช้เคลื่อนประจุ 1 คูลอมบ์ จากจุดที่ไม่มีสนามไฟฟ้ามายังจุดนั้น
\item ข) พลังงานที่ใช้เคลื่อนประจุ 1 คูลอมบ์ จากจุดอ้างอิงมายังจุดนั้น
\item ค) แรงที่ใช้เคลื่อนประจุ 1 คูลอมบ์ จากจุดอ้างอิงมายังจุดนั้น
\item ง) กำลังที่ใช้เคลื่อนประจุ 1 คูลอมบ์ จากจุดอ้างอิงมายังจุดนั้น
\end{enumerate}
\item จุดประจุ Q = +4 μC สร้างศักย์ไฟฟ้าที่ระยะ 2 m จากประจุ ศักย์ไฟฟ้ามีค่าเท่าใด? (k = 9×10⁹ N·m²/C²) \hfill \textbf{\small (Main)}
\begin{enumerate}[label=)]
\item ก) 9×10³ V
\item ข) 1.8×10⁴ V
\item ค) 4.5×10³ V
\item ง) 3.6×10⁴ V
\end{enumerate}
\item ประจุบวก q = 2 μC เคลื่อนจากจุด A (ศักย์ 500 V) ไปยังจุด B (ศักย์ 200 V) ในสนามไฟฟ้าสม่ำเสมอ งานที่สนามไฟฟ้าทำต่อประจุนี้มีค่าเท่าใด? \hfill \textbf{\small (Main)}
\begin{enumerate}[label=)]
\item ก) 6×10⁻⁴ J
\item ข) 6×10⁻³ J
\item ค) 6×10⁻² J
\item ง) 1.4×10⁻³ J
\end{enumerate}
\item ความจุไฟฟ้า (Capacitance) ของตัวเก็บประจุนิ่ม หมายถึงข้อใด? \hfill \textbf{\small (Main)}
\begin{enumerate}[label=)]
\item ก) ความสามารถในการเก็บพลังงาน
\item ข) ความสามารถในการเก็บประจุต่อหน่วยความต่างศักย์
\item ค) ความสามารถในการกั้นกระแสไฟฟ้า
\item ง) ความสามารถในการเคลื่อนประจุ
\end{enumerate}
\item ตัวเก็บประจุมีความจุ 10 μF ต่อกับความต่างศักย์ 50 V จะเก็บประจุได้กี่คูลอมบ์? \hfill \textbf{\small (Main)}
\begin{enumerate}[label=)]
\item ก) 5×10⁻⁴ C
\item ข) 5×10⁻³ C
\item ค) 5×10⁻² C
\item ง) 5 C
\end{enumerate}
\item ตัวเก็บประจุแผ่นขนานมีพื้นที่แผ่น 0.01 m² วางห่างกัน 0.001 m มีไดอิเล็กทริกเป็นอากาศ (ε₀ = 8.85×10⁻¹² F/m) ความจุของตัวเก็บประจุนี้เท่ากับเท่าใด? \hfill \textbf{\small (Main)}
\begin{enumerate}[label=)]
\item ก) 8.85×10⁻¹¹ F
\item ข) 8.85×10⁻¹⁰ F
\item ค) 8.85×10⁻¹² F
\item ง) 8.85×10⁻⁹ F
\end{enumerate}
\item ประจุสองประจุชนิดเดียวกันวางห่างกัน r ถ้าระยะห่างเพิ่มเป็น 2r แรงระหว่างประจุจะเปลี่ยนแปลงอย่างไร? \hfill \textbf{\small (Main)}
\begin{enumerate}[label=)]
\item ก) ลดลง 2 เท่า
\item ข) ลดลง 4 เท่า
\item ค) เพิ่มขึ้น 2 เท่า
\item ง) เพิ่มขึ้น 4 เท่า
\end{enumerate}
\item ข้อใดแสดงทิศทางของสนามไฟฟ้าที่เกิดจากประจุบวกได้ถูกต้อง? \hfill \textbf{\small (Main)}
\begin{enumerate}[label=)]
\item ก) ชี้ออกจากประจุ
\item ข) ชี้เข้าหาประจุ
\item ค) วงรอบประจุ
\item ง) ไม่มีทิศทาง
\end{enumerate}
\item ที่ระยะ r₁ = 1 m และ r₂ = 2 m จากประจุจุด Q ศักย์ไฟฟ้าต่างกัน 4500 V ค่าของ Q เป็นเท่าใด? (k = 9×10⁹) \hfill \textbf{\small (Main)}
\begin{enumerate}[label=)]
\item ก) 1 μC
\item ข) 2 μC
\item ค) 3 μC
\item ง) 4 μC
\end{enumerate}
\item ประจุสองประจุ +q และ -q วางห่างกัน 2d ณ จุดกึ่งกลางระหว่างประจุทั้งสอง ขนาดสนามไฟฟ้ารวมเป็นเท่าใด? \hfill \textbf{\small (Main)}
\begin{enumerate}[label=)]
\item ก) 0
\item ข) 2kq/d²
\item ค) 4kq/d²
\item ง) kq/d²
\end{enumerate}
\item ประจุสามประจุ +q, +q และ -q วางที่จุดยอดของสามเหลี่ยมด้านเท่า ขนาดของแรงลัพธ์ที่กระทำต่อประจุ -q เป็นเท่าใด ให้แรงระหว่างประจุที่เท่ากันแต่ละคู่คือ F₀ \hfill \textbf{\small (Main)}
\begin{enumerate}[label=)]
\item ก) F₀
\item ข) √3·F₀
\item ค) 2F₀
\item ง) 3F₀
\end{enumerate}
\item ตัวเก็บประจุแผ่นขนานต่อกับแบตเตอรี่ 100 V เมื่อเพิ่มระยะห่างระหว่างแผ่นเป็น 2 เท่า โดยไม่ตัดแบตเตอรี่ออก ประจุที่เก็บจะเปลี่ยนแปลงอย่างไร? \hfill \textbf{\small (Main)}
\begin{enumerate}[label=)]
\item ก) เพิ่มเป็น 2 เท่า
\item ข) เท่าเดิม
\item ค) ลดลง 2 เท่า
\item ง) ลดลง 4 เท่า
\end{enumerate}
\item ศักย์ไฟฟ้าที่ระยะอนันต์จากประจุจุดกำหนดให้เท่ากับศูนย์ ณ ตำแหน่งใดที่ศักย์ไฟฟ้ามีค่าเป็นศูนย์ ระหว่างประจุ +2 μC และ +8 μC ที่วางห่างกัน 90 cm? \hfill \textbf{\small (Main)}
\begin{enumerate}[label=)]
\item ก) 15 cm จากประจุ +2 μC
\item ข) 30 cm จากประจุ +2 μC
\item ค) 45 cm จากประจุ +2 μC
\item ง) 60 cm จากประจุ +2 μC
\end{enumerate}
\item เมื่อนำแท่งแก้วไปขัดด้วยผ้าไหม แท่งแก้วจะมีประจุชนิดใด? \hfill \textbf{\small (Main)}
\begin{enumerate}[label=)]
\item ก) ประจุลบ
\item ข) ประจุบวก
\item ค) เป็นกลาง
\item ง) ขึ้นอยู่กับชนิดของผ้า
\end{enumerate}
\item เส้นสนามไฟฟ้าที่ออกจากประจุบวกจะไปสิ้นสุดที่ใด? \hfill \textbf{\small (Main)}
\begin{enumerate}[label=)]
\item ก) ที่อนันต์
\item ข) ที่ประจุลบ
\item ค) ที่จุดศูนย์กลางของประจุ
\item ง) ที่ผิวของตัวนำ
\end{enumerate}
\item ตัวนำทรงกลมรัศมี R มีประจุ +Q กระจายอยู่บนผิว ณ จุดที่ห่างจากจุดศูนย์กลาง 0.5R สนามไฟฟ้ามีค่าเท่าใด? \hfill \textbf{\small (Main)}
\begin{enumerate}[label=)]
\item ก) kQ/R²
\item ข) 2kQ/R²
\item ค) 4kQ/R²
\item ง) 0
\end{enumerate}
\item ตัวนำทรงกลมรัศมี R มีประจุ +Q กระจายอยู่บนผิว ศักย์ไฟฟ้าภายในตัวนำ (r < R) มีค่าเท่าใด? \hfill \textbf{\small (Main)}
\begin{enumerate}[label=)]
\item ก) kQ/R
\item ข) kQ/r
\item ค) 0
\item ง) kQ/R²
\end{enumerate}
\item ประจุสองประจุ +Q และ +4Q วางห่างกัน 3 m ประจุที่สามต้องวางที่ใดเพื่อให้แรงลัพธ์บนประจุ +Q เป็นศูนย์? \hfill \textbf{\small (Main)}
\begin{enumerate}[label=)]
\item ก) ห่างจาก +Q ไปทาง +4Q 1 m
\item ข) ห่างจาก +Q ไปทางตรงข้าม +4Q 1 m
\item ค) ห่างจาก +4Q ไปทาง +Q 1 m
\item ง) ห่างจาก +Q ไปทาง +4Q 2 m
\end{enumerate}
\item ประจุไฟฟ้ามูลฐาน (Elementary charge) มีขนาดเท่ากับเท่าใด? \hfill \textbf{\small (Extra)}
\begin{enumerate}[label=)]
\item ก) 1.6×10⁻¹⁹ C
\item ข) 9.1×10⁻³¹ C
\item ค) 1.67×10⁻²⁷ C
\item ง) 8.85×10⁻¹² C
\end{enumerate}
\item เมื่อนำแท่งอำพัน (น้ำยาง) ไปขัดด้วยผ้าขนสัตว์ แท่งอำพันจะมีประจุชนิดใด? \hfill \textbf{\small (Extra)}
\begin{enumerate}[label=)]
\item ก) ประจุบวก
\item ข) ประจุลบ
\item ค) ประจุเป็นกลาง
\item ง) ประจุทั้งบวกและลบ
\end{enumerate}
\item ค่าคงตัว k ในกฎคูลอมบ์ในหน่วย SI มีค่าเท่าใด? \hfill \textbf{\small (Extra)}
\begin{enumerate}[label=)]
\item ก) 9×10⁸ N·m²/C²
\item ข) 9×10⁹ N·m²/C²
\item ค) 6.67×10⁻¹¹ N·m²/C²
\item ง) 8.85×10⁻¹² N·m²/C²
\end{enumerate}
\item ถ้าประจุทั้งสองเพิ่มเป็น 3 เท่า และระยะห่างเพิ่มเป็น 2 เท่า แรงระหว่างประจุจะเปลี่ยนเป็นกี่เท่าของเดิม? \hfill \textbf{\small (Extra)}
\begin{enumerate}[label=)]
\item ก) 1.5 เท่า
\item ข) 2.25 เท่า
\item ค) 4.5 เท่า
\item ง) 9 เท่า
\end{enumerate}
\item สนามไฟฟ้าสม่ำเสมอ (Uniform Electric Field) หมายถึงสนามที่มีลักษณะอย่างไร? \hfill \textbf{\small (Extra)}
\begin{enumerate}[label=)]
\item ก) สนามที่มีขนาดและทิศคงที่ตลอด
\item ข) สนามที่มีขนาดคงที่แต่ทิศเปลี่ยน
\item ค) สนามที่มีทิศคงที่แต่ขนาดเปลี่ยน
\item ง) สนามที่เปลี่ยนแปลงตามระยะทาง
\end{enumerate}
\item อิเล็กตรอน (ประจุ -e, มวล mₑ) เคลื่อนที่ในสนามไฟฟ้าสม่ำเสมอ E จากจุดหยุดนิ่ง อิเล็กตรอนจะเคลื่อนที่ไปทางทิศใด? \hfill \textbf{\small (Extra)}
\begin{enumerate}[label=)]
\item ก) ตามทิศสนามไฟฟ้า
\item ข) ตรงข้ามทิศสนามไฟฟ้า
\item ค) ตั้งฉากกับทิศสนามไฟฟ้า
\item ง) ไม่เคลื่อนที่
\end{enumerate}
\item ประจุบวก q = +3 μC เคลื่อนที่ในสนามไฟฟ้าสม่ำเสมอ E = 2000 N/C ตามทิศสนามไฟฟ้าเป็นระยะ 0.1 m งานที่สนามไฟฟ้าทำเป็นเท่าใด? \hfill \textbf{\small (Extra)}
\begin{enumerate}[label=)]
\item ก) 6×10⁻⁴ J
\item ข) 6×10⁻³ J
\item ค) 6×10⁻² J
\item ง) 6 J
\end{enumerate}
\item ประจุจุด Q สร้างศักย์ไฟฟ้า 1800 V ที่ระยะ 0.5 m และ 450 V ที่ระยะ x จงหาค่า x \hfill \textbf{\small (Extra)}
\begin{enumerate}[label=)]
\item ก) 1 m
\item ข) 2 m
\item ค) 3 m
\item ง) 4 m
\end{enumerate}
\item ความสัมพันธ์ระหว่างสนามไฟฟ้า (E) และศักย์ไฟฟ้า (V) ในกรณีสนามสม่ำเสมอคือข้อใด? \hfill \textbf{\small (Extra)}
\begin{enumerate}[label=)]
\item ก) E = V·d
\item ข) E = V/d
\item ค) E = V·q
\item ง) E = V²/q
\end{enumerate}
\item Farad (F) เทียบเท่ากับหน่วยใดในหน่วยฐาน SI? \hfill \textbf{\small (Extra)}
\begin{enumerate}[label=)]
\item ก) C·V
\item ข) C/V
\item ค) V/C
\item ง) J/C
\end{enumerate}
\item ตัวเก็บประจุแผ่นขนานมีความจุ C เมื่อเพิ่มพื้นที่แผ่นเป็น 3 เท่า โดยระยะห่างคงเดิม ความจุจะเปลี่ยนเป็นเท่าใด? \hfill \textbf{\small (Extra)}
\begin{enumerate}[label=)]
\item ก) C/3
\item ข) C
\item ค) 3C
\item ง) 9C
\end{enumerate}
\item ตัวเก็บประจุแผ่นขนานมีความจุ C เมื่อเพิ่มระยะห่างเป็น 3 เท่า โดยพื้นที่คงเดิม ความจุจะเปลี่ยนเป็นเท่าใด? \hfill \textbf{\small (Extra)}
\begin{enumerate}[label=)]
\item ก) C/3
\item ข) C
\item ค) 3C
\item ง) 9C
\end{enumerate}
\item ประจุสามประจุ +q, +q และ +q วางที่จุดยอดสามเหลี่ยมด้านเท่า ขนาดแรงลัพธ์ที่กระทำต่อประจุแต่ละตัวเป็นเท่าใด ให้แรงระหว่างประจุแต่ละคู่คือ F₀ \hfill \textbf{\small (Extra)}
\begin{enumerate}[label=)]
\item ก) F₀
\item ข) √3·F₀
\item ค) 2F₀
\item ง) 0
\end{enumerate}
\item ประจุ +Q วางที่จุดกำเนิด และประจุ -2Q วางที่ x = d บนแกน x ที่ตำแหน่งใดบนแกน x สนามไฟฟ้ามีค่าเป็นศูนย์? \hfill \textbf{\small (Extra)}
\begin{enumerate}[label=)]
\item ก) x < 0
\item ข) 0 < x < d
\item ค) x > d
\item ง) ไม่มีตำแหน่งใด
\end{enumerate}
\item ประจุ +Q และ -Q วางห่างกัน 2d บนแกน x ที่จุด x = d (กึ่งกลาง) ศักย์ไฟฟ้ารวมเป็นเท่าใด? \hfill \textbf{\small (Extra)}
\begin{enumerate}[label=)]
\item ก) kQ/d
\item ข) 2kQ/d
\item ค) 0
\item ง) 4kQ/d
\end{enumerate}
\item ตัวเก็บประจุแผ่นขนานต่อกับแบตเตอรี่ V เมื่อถอดแบตเตอรี่ออก แล้วเพิ่มระยะห่างระหว่างแผ่นเป็น 2 เท่า ความต่างศักย์ระหว่างแผ่นจะเปลี่ยนเป็นเท่าใด? \hfill \textbf{\small (Extra)}
\begin{enumerate}[label=)]
\item ก) V/2
\item ข) V
\item ค) 2V
\item ง) 4V
\end{enumerate}
\item ประจุไฟฟ้าที่สะสมบนตัวนำจะอยู่ที่ตำแหน่งใด? \hfill \textbf{\small (Extra)}
\begin{enumerate}[label=)]
\item ก) กระจายทั่วทั้งตัวนำ
\item ข) กระจายอยู่ภายในตัวนำ
\item ค) กระจายอยู่บนผิวนอกของตัวนำ
\item ง) กระจายอยู่ที่จุดศูนย์กลาง
\end{enumerate}
\item ทรงกลมโลหะ A รัศมี R มีประจุ +Q ถูกเชื่อมต่อด้วยสายไฟกับทรงกลมโลหะ B รัศมี 2R ที่ไม่มีประจุ หลังเชื่อมประจุจะกระจายอย่างไร? \hfill \textbf{\small (Extra)}
\begin{enumerate}[label=)]
\item ก) A ได้ Q/2, B ได้ Q/2
\item ข) A ได้ Q/3, B ได้ 2Q/3
\item ค) A ได้ 2Q/3, B ได้ Q/3
\item ง) ทั้งหมดอยู่ที่ A
\end{enumerate}
\item ประจุสี่ประจุ +q, +q, -q, -q วางที่จุดยอดสี่เหลี่ยมจัตุรัมด้าน a จุดศูนย์กลางของสี่เหลี่ยม สนามไฟฟ้ามีขนาดเท่าใด? \hfill \textbf{\small (Extra)}
\begin{enumerate}[label=)]
\item ก) 0
\item ข) 4kq/a²
\item ค) 4√2·kq/a²
\item ง) 8kq/a²
\end{enumerate}
\item ตัวเก็บประจุ 10 μF ต่อขนานกับตัวเก็บประจุ 20 μF ความจุรวมเป็นเท่าใด? \hfill \textbf{\small (Extra)}
\begin{enumerate}[label=)]
\item ก) 10 μF
\item ข) 20 μF
\item ค) 30 μF
\item ง) 6.67 μF
\end{enumerate}
\item ตัวเก็บประจุแผ่นขนานมีพื้นที่ A และระยะห่าง d ต่อกับแบตเตอรี่ V แรงที่ดึงแผ่นทั้งสองเข้าหากันมีค่าเท่าใด? \hfill \textbf{\small (Extra)}
\begin{enumerate}[label=)]
\item ก) ε₀AV²/(2d²)
\item ข) ε₀AV²/d²
\item ค) ε₀A²V/(2d)
\item ง) ε₀AV/d²
\end{enumerate}
\item ข้อใดไม่ใช่คุณสมบัติของเส้นสนามไฟฟ้า? \hfill \textbf{\small (Extra)}
\begin{enumerate}[label=)]
\item ก) ออกจากประจุบวก เข้าสู่ประจุลบ
\item ข) เส้นสนามไฟฟ้าตัดกันได้
\item ค) เส้นสนามไฟฟ้าทำมุมฉากกับพื้นผิวตัวนำ
\item ง) ความหนาแน่นของเส้นสนามแสดงขนาดของสนาม
\end{enumerate}
\item ตัวเก็บประจุแผ่นขนานระยะห่าง d มีประจุ Q ตัวเก็บประจุนี้ถูกเพิ่มระยะห่างเป็น 2d โดยไม่ให้ประจุรั่วไหล งานที่ทำโดยสนามไฟฟ้าเป็นเท่าใด? \hfill \textbf{\small (Extra)}
\begin{enumerate}[label=)]
\item ก) Q²d/(2ε₀A)
\item ข) Q²d/(ε₀A)
\item ค) Q²d/(4ε₀A)
\item ง) 0
\end{enumerate}
\item ประจุ +q และ -2q วางห่างกัน L แรงที่ประจุ -2q กระทำต่อประจุ +q มีขนาดเป็นกี่เท่าของแรงที่ +q กระทำต่อ -2q? \hfill \textbf{\small (Extra)}
\begin{enumerate}[label=)]
\item ก) 0.5 เท่า
\item ข) 1 เท่า
\item ค) 2 เท่า
\item ง) 4 เท่า
\end{enumerate}
\item โพรทอน (ประจุ +e) เคลื่อนที่ด้วยความเร็วต้น v₀ เข้าสู่สนามไฟฟ้าสม่ำเสมอ E ในทิศตรงข้ามกับสนาม โพรทอนจะเคลื่อนที่อย่างไร? \hfill \textbf{\small (Extra)}
\begin{enumerate}[label=)]
\item ก) ความเร็วคงที่
\item ข) ช้าลงอย่างสม่ำเสมอ
\item ค) เร็วขึ้นอย่างสม่ำเสมอ
\item ง) เคลื่อนที่เป็นวงกลม
\end{enumerate}
\item ตัวเก็บประจุแผ่นขนานมีไดอิเล็กทริก (εᵣ = 4) อัดเข้าไปเต็มพื้นที่ระหว่างแผ่น ขณะที่ยังต่อกับแบตเตอรี่ V ความจุจะเปลี่ยนเป็นเท่าใด? \hfill \textbf{\small (Extra)}
\begin{enumerate}[label=)]
\item ก) C/4
\item ข) C
\item ค) 4C
\item ง) 2C
\end{enumerate}
\item ประจุ Q กระจายอย่างสม่ำเสมอบนทรงกลมรัศมี R ณ ตำแหน่ง r = 2R (นอกทรงกลม) สนามไฟฟ้าและศักย์ไฟฟ้าเป็นเท่าใด? \hfill \textbf{\small (Extra)}
\begin{enumerate}[label=)]
\item ก) E = kQ/R², V = kQ/R
\item ข) E = kQ/(4R²), V = kQ/(2R)
\item ค) E = kQ/(4R²), V = kQ/(2R)
\item ง) E = kQ/(2R²), V = kQ/(2R)
\end{enumerate}
\item ประจุสองประจุ +q และ +4q วางห่างกัน L ต้องวางประจุที่สาม q' ที่ใดเพื่อให้แรงลัพธ์บน +q เป็นศูนย์? \hfill \textbf{\small (Extra)}
\begin{enumerate}[label=)]
\item ก) ระหว่าง +q กับ +4q ห่างจาก +q เป็น L/4
\item ข) ระหว่าง +q กับ +4q ห่างจาก +q เป็น L/3
\item ค) นอกช่วงระหว่าง +q กับ +4q ด้านของ +q
\item ง) นอกช่วงระหว่าง +q กับ +4q ด้านของ +4q
\end{enumerate}
\item ไดโพลไฟฟ้า (Electric Dipole) ประกอบด้วยประจุ +q และ -q ห่างกัน d ที่จุดกึ่งกลางบนแกนกลาง สนามไฟฟ้ามีขนาดเท่าใด? \hfill \textbf{\small (Extra)}
\begin{enumerate}[label=)]
\item ก) kq/d²
\item ข) 2kq/d²
\item ค) kq/(d²/4)
\item ง) 0
\end{enumerate}
\item ตัวเก็บประจุแผ่นขนานต่อกับแบตเตอรี่ 100 V มีประจุ 10⁻⁴ C ถ้าถอดแบตเตอรี่ออกแล้วเพิ่มระยะห่างเป็น 3 เท่า ความต่างศักย์ระหว่างแผ่นจะเป็นเท่าใด? \hfill \textbf{\small (Extra)}
\begin{enumerate}[label=)]
\item ก) 100 V
\item ข) 200 V
\item ค) 300 V
\item ง) 400 V
\end{enumerate}
\end{enumerate}
\end{document}
